% =====================================================================
\chapter{El experimento: VAR y causalidad de Granger}
\chaptermark{VAR y causalidad de Granger}
\label{cap:var}

Con los cuatro instrumentos construidos, la pregunta de la tesina se puede contrastar: ¿el tono
de las noticias anticipa el precio y el volumen de las acciones, o es al revés? Este capítulo
une los dos lados, prepara las series, presenta el modelo y sus resultados, y resume las pruebas
complementarias y la robustez.

% ---------------------------------------------------------------------
\section{La unión de los instrumentos}
\label{sec:union}

Los dos lados se unieron en la ventana en la que ambos tienen datos: 244 ruedas, del
05/05/2025 al 30/04/2026. Todo quedó en dos tablas: una por rueda (el mercado) y otra por acción
y rueda (11.468 filas). Del lado de las noticias, las variables son el tono anormal, la atención
anormal y el tono propio de las 17 empresas (capítulo~\ref{cap:finbeto}). Del lado del mercado,
el retorno diario del S\&P~Merval y de cada acción, y el \textit{volumen anormal}:
$\log(1+\text{volumen}_t)$ menos su promedio en las 20 ruedas previas, para que sea comparable
entre acciones y no crezca con la inflación. El régimen del capítulo~\ref{cap:hmm} marca cada
rueda, y se marcaron aparte los dos días electorales, su eco y abril de 2026. En la ventana, el
retorno del S\&P~Merval promedió 0,16\% diario con un desvío de 2,79\%, y sus extremos son los
dos días electorales ($-13{,}25\%$ y $+21{,}77\%$).

\textbf{Estacionariedad.} Un VAR necesita series estacionarias, que oscilen alrededor de un nivel
estable: si no, dos variables con tendencia parecen relacionadas aunque no lo estén. Se usaron dos
pruebas que preguntan al revés: la ADF \citep{dickey1979}, cuya hipótesis nula es que la serie no
es estacionaria, y la KPSS \citep{kpss1992}, cuya nula es que sí lo es. En el lado bursátil, con
tres años de datos, 46 de los 47 precios no son estacionarios y los 47 retornos sí: por eso los
modelos usan retornos y no precios \citep{zablotsky2002}. En la ventana de la unión se probaron
156 series: retornos, volúmenes anormales, atención y tono propio son estables. El tono en nivel
dio \textit{inconcluso}, porque cambia en ciclos de meses (promedio trimestral de $+7{,}5$, $+0{,}5$,
$+7{,}6$, $+3{,}7$ y $-2{,}9$), un rasgo conocido de las series de tono \citep{correarocha2026}; por
eso el modelo usa el \textit{tono anormal}, que es estable sin discusión. El VIX también dio
inconcluso y entra a los modelos solo a través del régimen.

% ---------------------------------------------------------------------
\section{El modelo}

Un VAR con $p$ rezagos explica el vector de variables de hoy, $y_t$, con sus valores de las $p$
ruedas anteriores:
\begin{equation}
y_t = c + A_1\, y_{t-1} + A_2\, y_{t-2} + \dots + A_p\, y_{t-p} + \varepsilon_t .
\end{equation}
Cuántas ruedas hacia atrás mira el modelo (los \textit{rezagos}) se eligió automáticamente entre
1 y 5 con el criterio de Akaike \citep{lutkepohl2005}, que busca el equilibrio entre ajustar bien
los datos y no complicar el modelo de más.

La prueba de Granger \citep{granger1969} responde una pregunta concreta: ¿conocer el tono de las
ruedas anteriores mejora la predicción del retorno de hoy, más allá de lo que ya predicen los
retornos anteriores? Si la mejora es estadísticamente significativa, el tono \textit{anticipa} al
retorno. La misma pregunta se hace al revés, para ver si el retorno anticipa al tono. La función
de impulso-respuesta, además, muestra cómo responde cada variable, rueda a rueda, a un golpe en
otra.

Se usaron 224 ruedas, del 02/06/2025 al 30/04/2026: las primeras 20 de la ventana quedan afuera
porque no tienen 20 ruedas previas con las que calcular el tono anormal. La prueba se repitió en
cinco niveles —el mercado, cada una de las 47 empresas, las 17 con tono propio, los 10 sectores y
los regímenes—, en las dos direcciones y con precio y volumen: 304 pruebas en total. Todos los
modelos resultaron estables, es decir, el efecto de un golpe se apaga con el tiempo en lugar de
crecer.

\textbf{Pruebas múltiples.} Con tantas pruebas, algunas salen significativas por azar: con 100
pruebas al 5\%, se esperan unas 5 aunque no haya ningún efecto \citep{harvey2016}. Por eso los
p-valores se corrigieron con Benjamini-Hochberg \citep{bh1995}, que los «encarece» para no
confundir casualidades con efectos; en todas las tablas, «significativo» quiere decir $p$
\textit{corregido} menor a 0,05. Como control, la prueba por empresa se repitió con el tono
mezclado al azar, lo que rompe cualquier relación real: el test debería indicar un efecto, por
error, en alrededor del 5\% de las pruebas, y dio 6,9\%. El método no encuentra efectos donde no
los hay.

% ---------------------------------------------------------------------
\section{¿Quién anticipa a quién?}
\label{sec:res-u6}

\subsection{El mercado}

De las ocho pruebas del VAR del mercado, una sola es significativa, y es fuerte: el retorno del
S\&P~Merval anticipa el tono anormal de las noticias ($p$ corregido 0,0003). Ninguna prueba en la
dirección de las noticias hacia el mercado llega a ser significativa (tabla~\ref{tab:u6-mercado}).

\begin{table}[H]
\centering\small
\begin{tabular}{@{}llccc@{}}
\toprule
\textbf{Prueba} & \textbf{Dirección} & \textbf{p} & \textbf{p corregido} & \textbf{Significativo} \\
\midrule
tono \flecha{} retorno & Noticias \flecha{} mercado & 0,2119 & 0,4238 & No \\
retorno \flecha{} tono & Mercado \flecha{} noticias & $<$0,0001 & 0,0003 & \textbf{Sí} \\
tono \flecha{} volumen & Noticias \flecha{} mercado & 0,0632 & 0,2528 & No \\
volumen \flecha{} tono & Mercado \flecha{} noticias & 0,3617 & 0,4823 & No \\
atención \flecha{} retorno & Noticias \flecha{} mercado & 0,9293 & 0,9293 & No \\
retorno \flecha{} atención & Mercado \flecha{} noticias & 0,1479 & 0,3944 & No \\
atención \flecha{} volumen & Noticias \flecha{} mercado & 0,3377 & 0,4823 & No \\
volumen \flecha{} atención & Mercado \flecha{} noticias & 0,5918 & 0,6764 & No \\
\bottomrule
\end{tabular}
\caption{Pruebas de causalidad de Granger en el VAR del mercado (224 ruedas). Tono: tono anormal
de las noticias de la rueda (tono menos el promedio de las 20 ruedas previas). Atención: cantidad
de notas comparada con lo habitual. Volumen: monto operado anormal. Fuente: elaboración propia
(\carpeta{Unión/06_var_granger}).}
\label{tab:u6-mercado}
\end{table}

La impulso-respuesta dimensiona el efecto (figura~\ref{fig:irf}). Un día en que el S\&P~Merval
sube un desvío (2,8\%), el tono anormal sube 1,84 puntos en la rueda siguiente y 1,23 en la
segunda, ambos distintos de cero; a la tercera ya no se distingue de cero. Como un desvío del
tono anormal son 6,0 puntos, la suba de la rueda siguiente es de un tercio de desvío. El tono y el
retorno de una misma rueda también se mueven juntos (correlación 0,19), pero con datos diarios no
se puede saber quién movió a quién dentro de la rueda. Desde la rueda siguiente, el efecto del
tono sobre el retorno ya no se distingue de cero.

\begin{figure}[H]
\centering
\includegraphics[width=0.8\linewidth]{u6_irf_mercado.png}
\caption{Funciones de impulso-respuesta del VAR del mercado ante un golpe de un desvío, con
banda del 95\%. Si la banda no toca el cero, la respuesta es distinta de cero. Fuente:
elaboración propia (\carpeta{Unión/06_var_granger}).}
\label{fig:irf}
\end{figure}

\subsection{Las empresas}

El retorno de casi todas las empresas anticipa el tono de las noticias: 39 de 47, todas con
signo positivo (tabla~\ref{tab:u6-empresas}). En este nivel el tono es el del mercado, y los
retornos de las empresas se mueven con el índice, de modo que en buena parte es el mismo hallazgo
del mercado repetido. En la dirección contraria, 7 de 94 pruebas dan $p<0{,}05$ sin corregir,
menos de las 9,4 que se esperarían por azar; después de la corrección quedan dos, sin un patrón.

\begin{table}[H]
\centering\small
\begin{tabular}{@{}lcccc@{}}
\toprule
\textbf{Prueba} & \textbf{Pruebas} & \textbf{$p<0{,}05$ sin corregir} & \textbf{Significativas} & \textbf{Signo} \\
\midrule
\multicolumn{5}{@{}l}{\textit{47 empresas, con el tono del mercado}} \\
tono \flecha{} retorno & 47 & 1 & 0 & — \\
tono \flecha{} volumen & 47 & 6 & 2 & 2 positivas \\
retorno \flecha{} tono & 47 & 43 & \textbf{39} & 39 positivas \\
volumen \flecha{} tono & 47 & 11 & 3 & 3 positivas \\
\midrule
\multicolumn{5}{@{}l}{\textit{17 empresas, con su tono propio}} \\
tono propio \flecha{} retorno & 17 & 0 & 0 & — \\
tono propio \flecha{} volumen & 17 & 2 & 1 & 1 positiva \\
retorno \flecha{} tono propio & 17 & 16 & \textbf{15} & 15 positivas \\
volumen \flecha{} tono propio & 17 & 1 & 0 & — \\
\bottomrule
\end{tabular}
\caption{Pruebas de Granger por empresa. «Significativas» cuenta las que siguen siéndolo
después de la corrección. Fuente: elaboración propia (\carpeta{Unión/06_var_granger}).}
\label{tab:u6-empresas}
\end{table}

La prueba más exigente es la del tono propio, que descuenta el clima general de las noticias. En
15 de las 17 empresas, cuando la acción sube, en las ruedas siguientes se habla de ella mejor que
del resto del mercado: BBAR, BMA, BYMA, CEPU, CRES, EDN, GGAL, IRSA, LOMA, METR, SUPV, TGSU2,
TRAN, TXAR e YPFD (no aparece en ALUA ni en PAMP). Es el resultado más fuerte de la tesina,
porque es específico de cada empresa (figura~\ref{fig:u6-empresas}). En la dirección contraria
queda una sola prueba significativa (IRSA, tono propio \flecha{} volumen), aislada.

\begin{figure}[H]
\centering
\includegraphics[width=0.42\linewidth]{u6_empresas.png}
\caption{Pruebas de Granger por empresa. Cada fila es una empresa y cada columna una prueba;
más oscuro, $p$ más chico; $\circ$ $p<0{,}05$ sin corregir; $\bullet$ significativo después
de la corrección. Fuente: elaboración propia (\carpeta{Unión/06_var_granger}).}
\label{fig:u6-empresas}
\end{figure}

\subsection{Los sectores y el régimen}

\textbf{Sectores.} Se repitió el VAR con el promedio de las acciones de cada uno de los diez
sectores. La dirección mercado \flecha{} noticias aparece en casi todos: en 9 de los 10, cuando el
sector sube, el tono de las ruedas siguientes es más positivo. La excepción es Tecnología, que
tiene una sola acción (BOLT). La dirección contraria aparece una sola vez: en Conglomerado el tono
anticipa el volumen, pero con signo negativo (más tono, menos volumen), un resultado aislado.

\textbf{Balance de todos los niveles.} Sumando mercado, empresas, tono propio y sectores, se
hicieron 152 pruebas en la dirección noticias \flecha{} mercado. Solo 4 son significativas, que es
lo que cabría esperar por azar con tantas pruebas, y no forman un patrón.

\textbf{Régimen.} La pregunta es si la relación cambia según el estado del mercado: la literatura
encuentra el efecto de las noticias sobre todo en los períodos de estrés \citep{garcia2013}. Con
los regímenes del capítulo~\ref{cap:hmm} se estimó un solo modelo con todas las ruedas y una marca
de día Tenso, para no pegar días que no son consecutivos; las 18 ruedas Moderado quedan afuera,
de modo que se comparan 166 días Calmo con 40 Tenso. Para cada prueba se pregunta tres cosas: si
hay efecto en los días Calmo, si hay efecto en los Tenso y si el efecto es distinto entre ambos.
La tabla~\ref{tab:u6-regimen} muestra los $p$ corregidos: un número chico indica que hay efecto, y
uno grande, que puede ser casualidad.

\begin{table}[H]
\centering\small
\begin{tabularx}{\linewidth}{@{}lcccY@{}}
\toprule
\textbf{Prueba} & \textbf{Días Calmo} & \textbf{Días Tenso} & \textbf{¿Cambia con el régimen?} & \textbf{Lectura} \\
\midrule
tono \flecha{} retorno & 0,42 & 0,92 & 0,95 & Nunca \\
tono \flecha{} volumen & 0,12 & 0,42 & 0,42 & Nunca \\
retorno \flecha{} tono & \textbf{0,007} & \textbf{0,007} & 0,81 & En los dos regímenes, igual \\
volumen \flecha{} tono & 0,60 & \textbf{0,008} & 0,12 & Solo en Tenso, sin diferencia significativa con Calmo \\
\bottomrule
\end{tabularx}
\caption{Pruebas de Granger según el régimen: $p$ corregido por Benjamini-Hochberg; en negrita,
menor a 0,05. «¿Cambia con el régimen?» prueba si el efecto es distinto en los días Tenso que en
los Calmo. Fuente: elaboración propia (\carpeta{Unión/06_var_granger}).}
\label{tab:u6-regimen}
\end{table}

La lectura es directa. \textbf{El tono no anticipa al mercado en ningún régimen}, ni siquiera en
los días de mayor estrés. \textbf{El retorno anticipa el tono en los dos regímenes y con la misma
fuerza}: los medios siguen al mercado tanto en calma como en tensión. El volumen anticipa el tono
solo en los días Tenso, pero la diferencia con los Calmo no llega a ser significativa: con 40 días
Tenso, los datos no alcanzan para afirmar que el régimen cambie la relación.

% ---------------------------------------------------------------------
\section{Dónde está el efecto: las notas de la misma noche}

La regla de las 17:00 manda las notas publicadas después del cierre a la rueda siguiente. Para
ver dónde está el efecto del retorno sobre el tono, se separaron las notas según la franja en que
se publicaron (tabla~\ref{tab:franjas}).

\begin{table}[H]
\centering\footnotesize
\begin{tabular}{@{}lrcc@{}}
\toprule
\textbf{Franja de publicación} & \textbf{Notas} & \textbf{Efecto de un desvío del retorno} & \textbf{p} \\
\midrule
Noche del día de la rueda (17 a 24 h) & 3.451 & \textbf{+4,77 puntos de tono} & $<$0,0001 \\
Día sin rueda & 2.142 & +3,11 & 0,1133 \\
Madrugada y mañana siguientes (antes de las 11) & 4.331 & +0,54 & 0,5142 \\
Sesión (11 a 17 h) & 4.001 & +1,26 & 0,1028 \\
\bottomrule
\end{tabular}
\caption{Relación entre el retorno de una rueda y el tono de las notas según la franja en que
se publicaron. Fuente: elaboración propia (\carpeta{Unión/08_robustez}).}
\label{tab:franjas}
\end{table}

El efecto está en las notas publicadas entre las 17:00 y la medianoche del día de la rueda: un
desvío del retorno sube su tono 4,8 puntos. En las notas de la madrugada y la mañana siguiente es
casi nulo. Por eso en el VAR aparece como «el retorno de hoy anticipa el tono de mañana»: después
del cierre, los medios cuentan cómo le fue al mercado y a cada empresa con un tono que acompaña la
suba o la baja, y no hay un arrastre que dure días.

% ---------------------------------------------------------------------
\section{Pruebas complementarias}
\label{sec:complementarias}

El VAR se complementó con tres pruebas (tabla~\ref{tab:complementarias}). La primera usa solo
las 8.069 notas publicadas con el mercado cerrado, que salen antes de cualquier operación de la
rueda: si su tono anticipa el salto de apertura, eso es noticias \flecha{} mercado sin ambigüedad
\citep{glasserman2025}. Las otras dos son estudios de eventos \citep{mackinlay1997}: alrededor
de las ruedas de tono extremo del mercado (16 eventos) y de cada empresa (88), con el retorno
normal estimado en las ruedas $[-25,-4]$ y el anormal medido antes, el día y después, contra
5.000 sorteos de fechas al azar; y alrededor de los shocks de cantidad con la definición del
plan ($z>2$ contra las 30 ruedas previas), medida por día de calendario porque la versión
original marcaba casi solo lunes. Dos de estas pruebas se especificaron por escrito antes de ver
los resultados.

\begin{table}[H]
\centering\footnotesize
\begin{tabularx}{\linewidth}{@{}L{3.3cm}YL{3.2cm}@{}}
\toprule
\textbf{Prueba} & \textbf{Resultado} & \textbf{Lectura} \\
\midrule
Notas con el mercado cerrado y retorno de apertura & Con notas de la noche más positivas el mercado abre algo más arriba (+0,18\% por desvío), pero no es significativo ($p=0{,}107$). Sin los dos días electorales, +0,06\% & El indicio se apoya en las elecciones: lo que movió la apertura fue el resultado, no el tono \\
Ruedas de tono extremo positivo (7 eventos) & El S\&P~Merval sube 2,01\% ese día (7 de 7; $p=0{,}018$ contra el azar) y ya venía subiendo 3,19\% en las 3 ruedas previas ($p=0{,}037$). Después, nada. Las negativas no muestran movimiento & El precio se mueve antes o junto con las noticias, no después \\
Tono propio extremo positivo (62 eventos) & La acción ya había subido 1,24\% más que el mercado en las 3 ruedas previas ($p=0{,}029$); después, $-$0,14\% & Se habla bien de la empresa después de que sube \\
Shock de cantidad en 17 empresas (105 eventos) & Volumen anormal del día 0,139 contra $-$0,054 de un día cualquiera ($p$ corregido 0,017); sin dirección en el precio, sin efecto después y sin efecto en la apertura & Apoyo parcial a H1 en el volumen: la cobertura acompaña a la actividad \\
\bottomrule
\end{tabularx}
\caption{Pruebas complementarias al VAR. Fuente: elaboración propia
(\carpeta{Unión/06b_noticias_preapertura}, \carpeta{07_event_study} y
\carpeta{07b_shock_volumen_noticias}).}
\label{tab:complementarias}
\end{table}

En los días electorales el caso es exactamente el que plantea la hipótesis: notas de la noche
muy negativas tras las bonaerenses (tono anormal $-5{,}3$, apertura $-6{,}6\%$) y muy positivas
tras las legislativas ($+8{,}5$, apertura $+13{,}4\%$). Pero lo que mueve al mercado es la
elección, no el tono con que los medios la cuentan; fuera de esos días, el tono de la noche se
relaciona con el cierre del día anterior (correlación 0,42) y no con la apertura. Lo mismo pasa en
el mercado con el shock de cantidad: su único salto de apertura depende del anuncio del Tesoro de
Estados Unidos del 24/09/2025. Un hecho grande genera muchas noticias y mueve el precio, pero lo
que mueve al mercado es el hecho, no la cobertura.

% ---------------------------------------------------------------------
\section{Robustez}
\label{sec:var-robustez}

La robustez cambió una sola decisión por vez, en doce variantes: tono ajustado por portal, tono en
nivel, sin notas de hora sospechosa, fecha de calendario en lugar de la regla de las 17:00, solo
notas del día antes del cierre, sin abril de 2026, días electorales como control, sin las acciones
ilíquidas, régimen filtrado, Calmo contra No calmo y dos umbrales para las ruedas de tono extremo.
De 106 comparaciones, 93 mantienen la conclusión. Dos resultados se sostienen
en todas: \textbf{en ninguna variante el tono anticipa el retorno}, ni en el mercado ni en ninguna
empresa, y el mercado anticipa al tono con cualquier medida de tono. Las únicas variantes en que
el efecto del retorno sobre el tono desaparece son las que dejan afuera, o cambian de rueda, las
notas publicadas después del cierre: la confirmación de que el efecto está en las notas de la
misma noche.


% ---------------------------------------------------------------------
\section{Discusión}

El resultado central coincide con uno de los dos hallazgos de \citet{tetlock2007}: los retornos
anticipan el tono de la cobertura. Que se repita con medios argentinos, en otro idioma y con un
modelo de sentimiento propio, lo hace más creíble; y la hora de publicación agrega una precisión
que una columna diaria no podía dar: el efecto está en las notas que cuentan la jornada. Es
consistente con \citet{verma2026}, para quienes buena parte de la cobertura refleja información
que el precio ya incorporó.

Que las noticias no anticipen al mercado difiere de \citet{tetlock2007} y de \citet{galdi2018},
pero no contradice a la literatura: el efecto se concentra en las recesiones \citep{garcia2013} y
la ventana es de un solo año. También es compatible con un mercado que incorpora la información
rápido \citep{fama1970}: la salida del cepo, las elecciones y el anuncio del Tesoro se reflejaron
de inmediato en el precio y en el VIX argentino. Que una empresa con muchas más notas se opere más
ese día está en línea con la literatura de atención \citep{da2011,engelberg2011}. Y el hallazgo por
empresa —después de que una acción sube, se habla de ella mejor que del resto— respalda medir el
sentimiento a nivel de entidad \citep{tang2023finentity}, algo que solo fue posible con el cruce
entre FinBETO y el NER. La corrección por pruebas múltiples y las especificaciones escritas de
antemano responden a la advertencia sobre la predictibilidad espuria \citep{nikolopoulos2026}.
